\section*{PHYSIQUE (13 points)}

\section{Propagation d'ondes mécaniques et d'ondes lumineuses}
Les ondes mécaniques et les ondes lumineuses sont caractérisées par des
propriétés bien déterminées. Les phénomènes liés à leur propagation permettent
de fournir des informations sur les milieux de propagation et la nature de la
lumière, et de déterminer certains paramètres caractéristiques.

Le but de l'exercice est de reconnaitre quelques propriétés des ondes
ultrasonores et des ondes lumineuses à partir de leur propagation dans
différents milieux.

\subsection{Propriétés des ondes ultrasonores et des ondes lumineuses}
Recopier sur votre copie, le numéro de la question, et écrire la lettre
correspondante à la seule proposition vraie parmi~:

\begin{table}[H]
    \centering
    \renewcommand{\arraystretch}{1.5}
    \begin{tabularx}{\textwidth}{|p{7mm}|X|}
        \hline
        $\mathbf{A}$ & Les ondes ultrasonores sont des ondes longitudinales. \\
        \hline
        $\mathbf{B}$ & Le domaine de fréquences de la lumière visible est limité
                       entre $\qty{400}{\nm}$ et $\qty{1000}{\nm}$. \\
        \hline
        $\mathbf{C}$ & Les ondes ultrasonores et les ondes lumineuses ont même
                       célérité de propagation dans le même milieu. \\
        \hline
        $\mathbf{D}$ & La fréquence des ondes lumineuses varie d'un milieu à un
                       autre. \\
        \hline
    \end{tabularx}
\end{table}



\subsection{Propagation des ondes ultrasonores}
On place en une même position, un émetteur $E$ et un récepteur $R$ des ondes
ultrasonores, à la distance $d=\qty{42,5}{\cm}$ d'un obstacle. Les ondes
ultrasonores qui se propagent à partir de $E$, se réfléchissent sur l'obstacle
puis sont reçues par $R$.

Un système d'acquisition informatique permet de visualiser l'onde émise~(a) et
l'onde reçue~(b). La figure~\ref{fig:oscillogramme_1} donne l'oscillogramme
obtenu.

\begin{enumerate}
    \item Déterminer la valeur du retard temporel $\tau$ entre les ondes~(a)
          et~(b).
    \item Vérifier que la valeur de la célérité de propagation dans l'air est
          $v_{\text{air}}=\qty{340}{\m\per\s}$.
    \item On répète l'expérience en utilisant le même dispositif, et l'eau comme
          milieu de propagation. On obtient avec le même système d'acquisition
          informatique l'oscillogramme représenté sur la
          figure~\ref{fig:oscillogramme_2}. Dans quel milieu (air/eau), la
          propagation des ondes ultrasonores est plus rapide~? Justifier votre
          réponse.
\end{enumerate}

\begin{minipage}{0.55\linewidth}
    \begin{figure}[H]
        \centering
        \includegraphics[width=.8\textwidth]{2025_10_31_838bdd6d133f8b78fc82g-3}
        \caption{}
        \label{fig:oscillogramme_1}
    \end{figure}
\end{minipage}
\begin{minipage}{0.43\linewidth}
    \begin{figure}[H]
        \centering
        \includegraphics[width=0.6\textwidth]{2025_10_31_838bdd6d133f8b78fc82g-3(1)}
        \caption{}
        \label{fig:oscillogramme_2}
    \end{figure}
\end{minipage}


\subsection{Propagation des ondes lumineuses}
On éclaire une fente verticale de largeur $a=\qty{0.1}{\mm}$, à l'aide d'un
laser qui donne une lumière monochromatique de longueur d'onde
$\lambda=\qty{632.8}{\nm}$. On observe sur un écran placé à la distance $D$ de
la fente, des taches lumineuses mettant en évidence le phénomène de diffraction.
La largeur de la tache centrale s'exprime par~: $\mathrm{L}=\frac{2\lambda
D}{a}$. La célérité de la lumière dans le vide (ou l'air) est
$c=\qty{3e8}{\m\per\s}$.

\begin{enumerate}
    \item Déterminer la valeur de la fréquence $v$ de la lumière utilisée.
    \item On refait l'expérience en utilisant un fil très fin vertical de
          diamètre $a_{0}$, on obtient une tache centrale de largeur $L_{0}=2L$.
          Déterminer la valeur de $a_{0}$.
\end{enumerate}




\section*{Exercice 2 ( $\mathbf{5}$ points) : réponse d'un dipôle}
L'étude électrique ou énergétique de quelques dipôles permet de déterminer certains paramètres qui les caractérisent, et de se rendre compte de leurs effets sur les phénomènes dont ces dipôles sont siège.\\
Le but de cet exercice est de déterminer l'inductance d'une bobine, et d'étudier la décharge d'un condensateur à travers cette bobine.

\section*{1. détermination de l'inductance d'une bobine}
Pour déterminer l'inductance $L$ d'une bobine de résistance négligeable, on utilise le montage représenté dans la figure (1), comprenant cette bobine, un conducteur ohmique de résistance $R=1,5 \cdot 10^{3} \Omega$, un GBF qui délivre une tension triangulaire de période $T$ et un interrupteur $K$. On ferme l'interrupteur $K$ à l'instant $t_{0}=0$, et on visualise à l'aide d'un oscilloscope la tension $u_{b}(t)$ aux bornes de la bobine, et la tension $u_{R}(t)$ aux bornes du conducteur ohmique. On obtient l'oscillogramme de la figure (2) (Page 5/7). - sensibilité verticale des deux voies de l'oscilloscope : $2 V . d i v^{-1}$.

\begin{itemize}
  \item balayage horizontale $0,2 \mathrm{~ms} . d \mathrm{v}^{-1}$
\end{itemize}

\begin{figure}[h]
\begin{center}
  \includegraphics[width=\textwidth]{2025_10_31_838bdd6d133f8b78fc82g-3(2)}
\captionsetup{labelformat=empty}
\caption{Figure 1}
\end{center}
\end{figure}

\subsection*{1.1. Quel est le rôle de la bobine lors de la fermeture du circuit?}
1.2. Montrer que les tensions $u_{R}$ et $u_{b}$ sont liées par la relation $u_{b}=-\frac{L}{R} \cdot \frac{d u_{R}}{d t}$.\\
1.3. Déterminer à partir de l'oscillogramme, les valeurs de $u_{b}$ et $\frac{d u_{R}}{d t}$ au cours de la première demipériode indiquée sur la figure (2).\\
1.4. Déduire que $L=0,1 H$.\\
2. Décharge d'un condensateur dans une bobine On réalise la décharge d'un condensateur dans la bobine précédente ( $L=0,1 H$ ) dans deux cas :\\
2.1. Premier cas : On utilise un condensateur de capacité $C$ initialement chargée sous la tension $U_{o}$, (fig.3). On note $q(t)$ la charge du condensateur à l'instant $t$.\\
2.1.1. Etablir l'équation différentielle vérifiée par la

\begin{figure}[h]
\begin{center}
  \includegraphics[width=\textwidth]{2025_10_31_838bdd6d133f8b78fc82g-4(2)}
\captionsetup{labelformat=empty}
\caption{Figure 2}
\end{center}
\end{figure}

2.1.2. Déterminer la valeur de $C$ sachant que le circuit est le siège d'oscillations électriques libres non amorties, de période propre $T_{0}=2 m s$. On prend $\pi^{2}=10$.\\
2.2. Deuxième cas : On utilise le condensateur précèdent de capacité $C$ initialement chargée sous la tension $U_{0}=6 \mathrm{~V}$, et on l'associe à la bobine précédente montée en série avec un conducteur ohmique de résistance $R$

\begin{figure}[h]
\begin{center}
  \includegraphics[width=\textwidth]{2025_10_31_838bdd6d133f8b78fc82g-4}
\captionsetup{labelformat=empty}
\caption{Figure 3}
\end{center}
\end{figure}

réglable et un interrupteur ouvert. On règle la résistance du conducteur ohmique sur la valeur $R_{0}$, et on ferme le circuit à l'instant $t_{0}=0$.\\
A l'aide d'un système d'acquisition informatique, on suit la tension $u_{C}(t)$ entre les bornes du condensateur, on obtient le graphe de la figure (4).\\
2.2.1. Nommer le régime d'oscillations que montre le graphe.\\
2.2.2. Calculer la valeur de l'énergie totale $\mathcal{E}_{0}$ du circuit à l'instant $t_{0}=0$ et la valeur de l'énergie totale $\mathcal{E}_{1}$ du circuit à l'instant $t_{1}=2 T$, avec $T$ pseudo période des oscillations électriques. Y at-il conservation de l'énergie totale du circuit ?\\
2.2.3. On admet que $\ln \left(\frac{\mathcal{E}_{0}}{\mathcal{C}_{1}}\right)=\frac{R_{0}}{L}\left(t_{1}-t_{0}\right)$.

Déterminer la valeur de $R_{0}$.

\begin{figure}[h]
\begin{center}
  \includegraphics[width=\textwidth]{2025_10_31_838bdd6d133f8b78fc82g-4(1)}
\captionsetup{labelformat=empty}
\caption{Figure 4}
\end{center}
\end{figure}

\section*{Exercice 3 ( $\mathbf{5}$ points) : Saut avec moto}
Le saut en longueur avec moto est considéré parmi les sports motivant, attirant et défiant pour dépasser certains obstacles naturels et artificiels.\\
Le but de cet exercice est d'étudier le mouvement du centre d'inertie $G$ d'un système ( $S$ )de masse $m$ constitué d'une moto avec motard sur une piste de course.\\
La piste de course est constituée d'une partie rectiligne horizontale, d'une partie rectiligne inclinée d'un angle $\alpha$ par rapport au plan horizontal et d'une zone de chute comportant un obstacle $(E)$ de hauteur $L$ situé à la distance $d$ de l'axe vertical passant par le point $D$, (fig1) (Page (6/7).

\section*{Données :}
\begin{itemize}
  \item Tous les frottements sont négligeables ;
\end{itemize}

$$
-\quad \alpha=26^{\circ} \quad ; \quad d=20 \mathrm{~m} \quad ; \quad L=10 \mathrm{~m} \quad ; \quad m=190 \mathrm{~kg}
$$

\begin{figure}[h]
\begin{center}
  \includegraphics[width=\textwidth]{2025_10_31_838bdd6d133f8b78fc82g-5(1)}
\captionsetup{labelformat=empty}
\caption{Figure 1}
\end{center}
\end{figure}

\section*{1. Mouvement du système ( $S$ ) sur la partie horizontale}
Le système ( $S$ ) démarre d'une position ou son centre d'inertie $G$ coïncide avec le point $A$. $G$ passe par le point $B$ avec la vitesse $\bar{v}_{0}=v_{0} . \bar{i}$ à l'instant $t_{0}=0$. Au cours de son mouvement, le système ( $S$ ) est soumis à une force motrice horizontale constante $\bar{F}$ ayant le même sens du mouvement. La trajectoire de $G$ est rectiligne.\\
Pour étudier le mouvement de $G$ entre $B$ et $C$ on choisit le repère (B.i) lié à la terre considéré comme galiléen. A $t_{0}=0$, on a : $x_{G}=x_{B}=0$.\\
1.1. En appliquant la deuxième loi de newton, montrer que l'expression de l'accélération de $G$ s'écrit : $a_{G}=\frac{F}{m}$. En déduire la nature du mouvement de $G$.\\
1.2. L'expression de la vitesse instantanée de $G$ s'écrit $v_{G}(t)=a_{G} . t+v_{0}$.\\
a. Choisir, en justifiant votre réponse, la courbe qui représente la vitesse instantanée $v_{G}(t)$ parmi les quatre courbes représentées sur la figure (2).

\begin{figure}[h]
\begin{center}
  \includegraphics[width=\textwidth]{2025_10_31_838bdd6d133f8b78fc82g-5}
\captionsetup{labelformat=empty}
\caption{Figure 2}
\end{center}
\end{figure}

b. En déduire les valeurs de la vitesse initiale $v_{0}$, et de l'accélération $a_{G}$ de $G$.\\
1.3. Calculer l'intensité de la force motrice $\bar{F}$.

\section*{2. Mouvement du système ( $S$ )durant la phase du saut}
Le système ( $S$ )quitte la piste de course au passage de $G$ par le point $D$ avec une vitesse $\dot{v_{D}}$ formant un angle $\alpha$ avec le plan horizontal pour sauter à travers l'obstacle ( $E$ ) (voir fig. (1)). Au cours du saut le système ( $S$ )n'est soumis qu'à son poids.

On étudie le mouvement de $G$ dans le champ de pesanteur uniforme dans un repère orthonormé (O.i.j) lié à la terre considéré comme galiléen. On choisit l'instant de passage de $G$ par le point $D$ comme nouvelle origine des dates $t_{0}=0$, tel que : $y_{0}=O D=h$.\\
2.1. En appliquant la deuxième loi de newton, montrer que les équations différentielles vérifiées par $x_{\mathrm{G}}(t)$ et $y_{\mathrm{G}}(t)$ coordonnées de $G$ dans le repère $(O . \dot{i . j})$ sont :

$$
\frac{d x_{G}}{d t}=v_{D} \cdot \cos \alpha \quad ; \quad \frac{d y_{G}}{d t}=-g \cdot t+v_{D} \cdot \sin \alpha
$$

2.2. L'expression numérique des équations horaires $x_{\mathrm{G}}(t)$ et $y_{\mathrm{G}}(t)$ du mouvement de $G$ est :

$$
x_{\mathrm{G}}(t)=22,5 . t(\mathrm{~m}) \quad ; \quad y_{\mathrm{G}}(t)=-5 . t^{2}+11 . t+5(\mathrm{~m})
$$

Déterminer les valeurs de la hauteur $h$, et de la vitesse $v_{D}$.\\
2.3. Le saut est réussi si la condition : $y_{G}>L+0,6(m)$ est vérifiée. Est-ce que le saut du motard est réussi ? Justifier votre réponse.



